\subsection{正則極限と可測包絡}
\label{sec:path-inputs}
\begin{proof}[命題\ref{input:paths}の証明]
\contractref{P1}について、全時間Cauchy性から厳密増加する $f$ を選び、
\[
 \mu\{\sup_t|Y_{f(n+1)}(t)-Y_{f(n)}(t)|>2^{-n}\}\leq2^{-n}
\]
とする。右連続性により上限は非負有理時刻の上限であり、事象は可測である。
Borel--Cantelliから、ほとんど全ての経路で、十分先の添字から隣接差の全時間上限が $2^{-n}$ 以下となる。
その尾部級数は有限なので一様Cauchy列となり、実数の完備性により全時間一様極限が存在する。
一様極限は右連続性と各正時刻の左極限を保つ。
共通の例外零集合は通常の条件により $\F_0$ に属するので、そこでは全時刻で零と定める。
各固定時刻で適合な実数値過程の極限であり、零集合上の修正も適合性を保つ。
これが強適合càdlàgな代表元を与える。

\contractref{P2}には同じ速い部分列を使う。各 $Y_n$ の経路は、コンパクト区間上のcàdlàg性と
無限時刻での有限極限により全時間有界である。
従って一様Cauchy性が成立する共通の確率1の集合上で
$\sup_{n,t}|Y_{f(n)}(t)|<\infty$ である。
右稠密な可算時刻集合 $D$ による拡張実数値可測関数
$\sup_{n,t\in D}|Y_{f(n)}(t)|$ を取り、有限な場合はその値、無限の場合は零とする。
この有限非負可測関数 $q$ は共通の確率1の集合上で全ての利得を支配する。

\contractref{P3}では $Z=E_\mu e^{-q}$、$dQ/d\mu=e^{-q}/Z$ と置く。
$q$ は有限なので $0<Z\leq1$、密度は正であり $Q\sim\mu$ となる。
$x^2e^{-x}$ は $[0,\infty)$ で有界なので $q\in L^2(Q)$ である。
同値測度は同じ零集合を持つため完備性を保ち、filtrationの右連続性は測度に依存しない。
従って通常の条件も保たれる。
\end{proof}
